% !TeX encoding = UTF-8
\documentclass[variant=apuntes,cover=institutional,mode=screen]{ua-engineering-v6-article}
% !TeX encoding = UTF-8
\UASetUniversity{Universidad Austral}
\UASetFaculty{Facultad de Ingeniería}
\UASetDepartment{Departamento de Computación}
\UASetCampus{Pilar}
\UASetCareer{Ingeniería Informática}
\UASetCommission{Comisión A}
\UASetProfessor{Equipo docente}
\UASetSemester{Segundo cuatrimestre 2026}
\UASetCourse{Sistemas Distribuidos}
\UASetDate{2026}

\UASetSemester{Segundo cuatrimestre 2026}
\UASetDate{2026}
\usetikzlibrary{positioning,arrows.meta,calc,fit,matrix,decorations.pathreplacing}

\UASetClassNumber{06}
\UASetClassTitle{Storage distribuido y almacenamiento permanente}
\UASetDocKind{Apuntes autocontenidos}
\title{Clase 06 --- Storage distribuido: almacenamiento permanente, acceso concurrente y rendimiento bajo fallas}
\author{\UAProfessor}
\date{\UADate}

\tikzset{
  uaN/.style={draw=UAIngNavy,rounded corners=1.5mm,thick,fill=white,
    minimum height=0.75cm,align=center,font=\small,inner sep=4pt},
  uaS/.style={uaN,fill=UAIngPaper},
  uaO/.style={uaN,draw=UAIngOrange,fill=UAIngPaper},
  uaA/.style={->,very thick,>=Latex,draw=UAIngNavy},
  uaAO/.style={->,very thick,>=Latex,draw=UAIngOrange}
}

\begin{document}
\setlength{\emergencystretch}{2em}
\maketitle
\tableofcontents
\clearpage

\section{Propósito y pregunta central}

Las primeras cinco clases del curso construyeron un lenguaje para razonar sobre incertidumbre, comunicación remota, reintentos, replicación, consenso y transacciones. Esta clase cambia deliberadamente de nivel: ya no pregunta solamente \emph{qué garantía queremos}, sino \emph{qué debe ocurrir en un nodo de almacenamiento para que esa garantía siga siendo verdadera después de una falla}.

Cuando una aplicación muestra ``compra confirmada'', cuando una réplica declara una entrada \emph{committed} o cuando una base de datos acepta un \texttt{COMMIT}, existe una cadena de mecanismos físicos y lógicos que debe transformar una intención en estado permanente, recuperable y consultable. En esa cadena aparecen buffers, registros, páginas, versiones, índices, archivos ordenados y tareas de mantenimiento.

\begin{uakeybox}
\textbf{Pregunta central.} ¿Cómo convierte un sistema distribuido sus operaciones en estado permanente, recuperable, accesible concurrentemente y eficiente bajo fallas?
\end{uakeybox}

La respuesta de esta clase no es un único producto ni una única estructura. Es una composición de responsabilidades:

\begin{enumerate}
\item \textbf{Registrar decisiones}: conservar evidencia suficiente para reconstruir una operación confirmada.
\item \textbf{Controlar acceso concurrente}: permitir que lectores y escritores avancen sin mezclar arbitrariamente estados.
\item \textbf{Elegir rutas físicas}: evitar revisar todos los datos cuando existe una estructura de acceso adecuada.
\item \textbf{Optimizar el camino de escritura}: transformar muchas escrituras pequeñas en trabajo físico sostenible.
\item \textbf{Proteger la capacidad}: aplicar mantenimiento y backpressure antes de que una cola interna crezca sin límite.
\end{enumerate}

Al finalizar, el estudiante debería poder explicar el problema, el mecanismo, la garantía, el costo y la evidencia para WAL, MVCC, índices, \texttt{EXPLAIN}, LSM-tree, compactación, filtros de Bloom y pausas de escritura.


\section{Cómo estudiar este capítulo}

Este capítulo no está pensado para memorizar una lista de siglas. Conviene estudiarlo en tres pasadas complementarias.

\begin{enumerate}
\item \textbf{Primera pasada: identificar la promesa observable.} Antes de mirar el mecanismo, formule qué debe seguir siendo verdadero para el usuario o para otro componente. En CampusShop, por ejemplo, una compra confirmada no debe desaparecer ni repetirse después de una caída.

\item \textbf{Segunda pasada: dibujar la ejecución y mover la falla.} Represente en un timeline qué ocurre en el cliente, el servicio, la memoria del motor, el WAL, las páginas y los procesos de mantenimiento. Desplace el crash, el timeout o la sobrecarga entre pasos y determine qué evidencia sobrevive.

\item \textbf{Tercera pasada: declarar frontera, costo y evidencia.} Para cada mecanismo responda: qué problema resuelve, qué garantía ofrece, dónde termina esa garantía, qué costo desplaza y qué métrica o experimento permitiría verificarla.
\end{enumerate}

Una lectura útil de cada sección sigue, por tanto, esta secuencia:

\[
\text{problema}\rightarrow\text{intuición}\rightarrow\text{regla técnica}
\rightarrow\text{falla}\rightarrow\text{evidencia}\rightarrow\text{decisión}.
\]

\begin{uainfo}
\textbf{Regla de estudio.} Si una explicación termina únicamente en ``usar WAL'', ``usar MVCC'', ``crear un índice'' o ``elegir una LSM-tree'', todavía falta justificar la garantía, la falla considerada y el costo físico u operacional que se acepta.
\end{uainfo}

\section{Mapa intuitivo antes del formalismo}

Antes de profundizar, es útil ubicar cada mecanismo dentro del problema general del almacenamiento distribuido. Las analogías siguientes sirven para iniciar la comprensión, pero cada una tiene un límite que luego formalizaremos.

\begin{description}
\item[WAL como cuaderno de evidencia.] El motor anota cómo reconstruir una decisión antes de depender de la escritura final de todas las páginas. La analogía no implica que el WAL sea un backup completo ni que el registro \texttt{COMMIT} contenga por sí solo todos los datos.

\item[MVCC como historial de ediciones.] En lugar de sobrescribir inmediatamente la versión anterior, el motor conserva varias versiones y aplica una regla de visibilidad a cada lector. No significa que se copie toda la base para cada transacción ni que desaparezcan todas las anomalías de concurrencia.

\item[Árbol B e índice como mapa; \texttt{EXPLAIN} como itinerario.] El índice reduce dónde buscar y el planificador elige una ruta física. \texttt{EXPLAIN} permite inspeccionar esa decisión, pero la existencia de un índice no obliga al motor a utilizarlo.

\item[LSM-tree como recepción rápida y organización posterior.] Las escrituras se acumulan en memoria y se vuelcan a archivos ordenados e inmutables; la compactación fusiona luego esos archivos. La escritura rápida no elimina trabajo: lo desplaza hacia lecturas, espacio y mantenimiento posterior.

\item[Backlog y write stall como control de admisión.] Si la compactación procesa menos trabajo del que ingresa, se acumula una cola interna. El motor puede retrasar o bloquear escritores para mantenerla acotada. La pausa protege al sistema, pero no crea capacidad nueva.
\end{description}

Este mapa permite formular cinco preguntas antes de introducir detalles de implementación:

\begin{enumerate}
\item ¿Qué evidencia debe sobrevivir para reconstruir una decisión confirmada?
\item ¿Qué versión puede observar cada lectura concurrente?
\item ¿Cómo se localizan las filas candidatas sin recorrer toda la tabla?
\item ¿Dónde y cuándo se paga el trabajo físico de organizar escrituras?
\item ¿Cómo se evita que el mantenimiento pendiente crezca sin límite?
\end{enumerate}

Las secciones siguientes convertirán estas intuiciones en reglas, timelines, estructuras y métricas verificables.

\section{Contexto: dónde vive el almacenamiento distribuido}

\subsection{Dos capas que no deben confundirse}

En un sistema distribuido conviven dos niveles de decisión.

\begin{center}
\begin{tikzpicture}[node distance=0.35cm and 0.55cm]
\node[uaO,text width=2.5cm] (op) {Operación lógica\\confirmar compra};
\node[uaS,text width=3.2cm,right=of op] (dist) {Capa distribuida\\particiona, replica, enruta y decide autoridad};
\node[uaS,text width=3.2cm,right=of dist] (local) {Motor local\\persiste, versiona, indexa y recupera};
\node[uaO,text width=2.6cm,right=of local] (state) {Estado permanente\\recuperable y consultable};
\draw[uaAO] (op)--(dist);
\draw[uaA] (dist)--(local);
\draw[uaAO] (local)--(state);
\end{tikzpicture}
\end{center}

\paragraph{Capa distribuida.} Decide dónde vive un dato, qué nodos lo replican, qué nodo o quorum tiene autoridad para confirmar y qué orden deben respetar las operaciones.

\paragraph{Motor de almacenamiento local.} Traduce decisiones lógicas en bytes persistentes, versiones visibles, páginas, índices, archivos y reglas de recuperación. WAL y MVCC son mecanismos del motor local; no constituyen por sí solos la distribución del sistema.

Esta separación evita dos errores frecuentes:

\begin{itemize}
\item creer que usar WAL convierte automáticamente una base local en una base distribuida;
\item creer que un protocolo de consenso vuelve seguro cualquier almacenamiento local aunque el nodo no persista correctamente votos, entradas o términos.
\end{itemize}

\subsection{Qué llamaremos motor de almacenamiento}

El \textbf{motor de almacenamiento} es el software que traduce lecturas y escrituras lógicas en:

\begin{itemize}
\item decisiones transaccionales;
\item bytes en memoria y en medio persistente;
\item registros de recuperación;
\item versiones y reglas de visibilidad;
\item páginas, árboles, índices y archivos ordenados;
\item procedimientos de mantenimiento y recuperación.
\end{itemize}

Un motor concreto puede ser parte de PostgreSQL, MySQL/InnoDB, RocksDB, Cassandra, una base embebida o un servicio de almacenamiento construido sobre otros componentes. Esta clase estudia principios transferibles y utiliza productos reales como ejemplos, no como recetas universales.

\subsection{Camino crítico y patrón de entrada/salida}

El \textbf{camino crítico de una escritura} es la secuencia de pasos que debe completarse antes de que el sistema pueda responder al cliente según la garantía prometida. No incluye necesariamente todo el trabajo físico asociado a la operación: parte de ese trabajo puede quedar diferido, siempre que exista evidencia durable suficiente para completarlo después.

Por ejemplo, el camino crítico de una compra puede incluir validar el invariante, registrar el cambio en WAL, volver durable el registro de COMMIT y enviar el ACK. La escritura posterior de todas las páginas e índices puede ocurrir fuera de ese camino.

También distinguiremos dos formas de entrada/salida:

\begin{itemize}
\item \textbf{E/S secuencial}: bytes contiguos se agregan o leen en orden, como al anexar registros a un log.
\item \textbf{E/S aleatoria o dispersa}: se accede a ubicaciones no contiguas, por ejemplo una página de datos, otra página distante y varias páginas de índices.
\end{itemize}

``Aleatoria'' no significa imprevisible. Describe la distribución física de los accesos. La diferencia de costo depende del dispositivo, la caché y la profundidad de colas, pero sigue siendo útil para entender por qué WAL y LSM desplazan trabajo hacia escrituras agrupables y secuenciales.

\section{Caso conductor: CampusShop}

CampusShop permite comprar almuerzos y entradas desde la aplicación del campus. La operación lógica se identifica con la clave estable \textbf{9F2A}. Al confirmar una compra, el sistema debe:

\begin{enumerate}
\item insertar una fila en \texttt{purchases};
\item descontar una unidad en \texttt{inventory};
\item actualizar un índice por estudiante, estado y fecha;
\item conservar un resultado asociado a 9F2A para reconocer reintentos;
\item responder ``compra confirmada'' solamente cuando la promesa sea recuperable.
\end{enumerate}

\begin{uainfo}
\textbf{Qué es 9F2A.} Es un identificador estable de la \emph{operación lógica}. Si el cliente pierde la respuesta y reintenta, vuelve a enviar 9F2A. No es una posición del WAL, una página física ni el contenido de una fila.
\end{uainfo}

Para discutir la implementación, usaremos dos etiquetas didácticas:

\begin{itemize}
\item \textbf{P17}: una página física que contiene parte de una tabla;
\item \textbf{I5}: una entrada de índice que ayuda a localizar P17.
\end{itemize}

La pregunta de apertura es deliberadamente sencilla:

\begin{quote}
La aplicación mostró ``compra confirmada''. Un segundo después se cortó la energía. ¿Qué debe seguir siendo verdadero al reiniciar?
\end{quote}

La respuesta de negocio es: la compra debe existir una sola vez, el stock no debe retroceder y el historial debe poder consultarse. Los mecanismos físicos aparecerán después de definir esa promesa.

\section{Vocabulario y siglas antes del detalle}

\subsection{Siglas centrales}

\begin{description}
\item[WAL --- Write-Ahead Logging.] En español, \textbf{registro anticipado de escritura}. El motor registra primero evidencia suficiente para reconstruir una modificación después de una falla.

\item[MVCC --- Multi-Version Concurrency Control.] En español, \textbf{control de concurrencia multiversión}. El motor conserva varias versiones de un dato y aplica reglas para decidir cuál puede observar cada transacción.

\item[LSM-tree --- Log-Structured Merge-tree.] En español, \textbf{árbol de fusión estructurado como log}. Recibe escrituras en memoria, las vuelca a archivos ordenados e inmutables y posteriormente fusiona esos archivos.

\item[ACK --- Acknowledgement.] Acuse de recibo. Es la respuesta que comunica al cliente que ya se cumplió la garantía prometida.

\item[LSN --- Log Sequence Number.] Número de secuencia que identifica una posición en el WAL. Ordena registros del log; no es un reloj físico.

\item[SQL --- Structured Query Language.] Lenguaje estructurado de consultas para leer y modificar datos relacionales.

\item[DDL --- Data Definition Language.] Parte de SQL utilizada para crear o modificar estructuras, por ejemplo \texttt{CREATE INDEX}.

\item[RPC --- Remote Procedure Call.] Llamada a procedimiento remoto. Una operación lógica atraviesa una red y puede terminar con respuesta, timeout o resultado desconocido.

\item[ACID.] Atomicidad, consistencia, aislamiento y durabilidad. Describe propiedades transaccionales; no explica por sí solo cómo se implementan físicamente.

\item[OLTP --- Online Transaction Processing.] Procesamiento transaccional en línea: muchas operaciones cortas, concurrentes y sensibles a latencia.

\item[SLO --- Service Level Objective.] Objetivo de nivel de servicio, por ejemplo ``99\% de escrituras por debajo de 20 ms''.

\item[RPO --- Recovery Point Objective.] Máxima pérdida de datos aceptable medida respecto de un punto de recuperación.

\item[RTO --- Recovery Time Objective.] Tiempo máximo aceptable para restaurar el servicio después de una falla.
\end{description}

\subsection{Confirmar y volver durable}

\paragraph{COMMIT.} Decisión del motor de aceptar definitivamente los efectos de una transacción. Un COMMIT puede estar decidido y ser durable antes de que todas las páginas de tabla e índice hayan sido escritas al archivo de datos.

\paragraph{X durable.} Decir que \emph{X es durable} significa que la evidencia necesaria para reconstruir X sobrevivirá a la falla declarada por el modelo. La palabra no tiene sentido sin responder:

\begin{enumerate}
\item ¿qué falla se considera: proceso, sistema operativo, energía, dispositivo, rack o región?;
\item ¿qué evidencia debe sobrevivir?;
\item ¿qué mecanismo demuestra que cruzó la frontera correspondiente?
\end{enumerate}

\paragraph{Log durable.} Prefijo del WAL cuyos bytes necesarios para recuperación sobrevivirían a la falla declarada.

\paragraph{Flush durable.} Volcado ya completado que llevó los bytes requeridos del WAL más allá de la frontera durable. Un flush común puede mover datos entre buffers sin hacerlos durables; la palabra \emph{durable} afirma que se alcanzó la capa que sobrevive a la falla considerada.

\paragraph{Frontera durable.} Punto después del cual la falla declarada ya no puede borrar la evidencia necesaria para reconstruir la operación. No es universal: depende del sistema operativo, \texttt{fsync}, controlador, dispositivo, configuración y modelo de falla.

\subsection{Escritura, buffers y recuperación}

\begin{description}
\item[\texttt{write}/append.] Copiar o agregar bytes a una capa. La llamada puede retornar mientras los datos continúan en memoria o caché.
\item[Flush o volcado.] Mover datos acumulados desde un buffer hacia una capa inferior. No implica necesariamente durabilidad.
\item[\texttt{fsync}.] Solicitar al sistema operativo que sincronice los datos relevantes con el medio persistente según la semántica del hardware y la configuración.
\item[Página sucia.] Página modificada en memoria que todavía no fue escrita al archivo de datos.
\item[Crash.] Caída abrupta de un proceso, del sistema operativo o de la energía. El modelo debe indicar qué capas se pierden.
\item[Recovery.] Recuperación: procedimiento de reinicio que usa la evidencia persistida para reconstruir un estado lógico válido.
\item[REDO.] Reaplicar una modificación confirmada que todavía no estaba materializada en su página o índice.
\item[UNDO.] Deshacer efectos de una transacción no confirmada cuando el diseño permite que esos efectos hayan llegado a páginas persistentes.
\item[Checkpoint.] Punto durable que permite acotar desde dónde debe comenzar la recuperación y cuánto WAL necesita procesarse.
\end{description}

\subsection{Concurrencia y consultas}

\begin{description}
\item[Versión.] Representación de una fila creada por una transacción y visible para ciertos lectores.
\item[Snapshot.] Instantánea lógica: regla que determina qué transacciones confirmadas y qué versiones puede observar una lectura. No es una copia física completa de la base.
\item[Vista coherente.] Resultado cuyas filas obedecen una misma regla de visibilidad, sin mezclar momentos arbitrarios.
\item[Serializabilidad.] Garantía de que el resultado concurrente equivale a alguna ejecución de las transacciones una detrás de otra.
\item[Planificador.] Componente que compara rutas físicas posibles y elige un plan estimado.
\item[\texttt{EXPLAIN}.] Muestra el plan elegido y sus estimaciones.
\item[\texttt{EXPLAIN ANALYZE}.] Ejecuta la consulta y agrega mediciones reales del recorrido.
\end{description}

\section{Repaso ágil: cinco puentes hacia el almacenamiento}

El repaso debe recuperar una sola idea por clase:

\begin{enumerate}
\item \textbf{Incertidumbre.} Si no sabemos si una operación ocurrió, el almacenamiento conserva evidencia para responder de manera estable.
\item \textbf{RPC y reintento.} Un reintento es seguro solo si identidad, efecto y resultado sobreviven a la falla.
\item \textbf{Replicación.} Varias copias no indican por sí solas cuáles son durables, actuales ni autorizadas.
\item \textbf{Consenso.} El protocolo decide qué debe recordarse; cada nodo debe persistirlo realmente.
\item \textbf{Transacciones.} ACID define un contrato; WAL y MVCC implementan parte de ese contrato.
\end{enumerate}

\begin{uakeybox}
Hoy dejamos de usar ``COMMIT'', ``durable'' y ``aislado'' como palabras abstractas y preguntamos qué bytes, esperas y estructuras las vuelven verdaderas.
\end{uakeybox}

\section{WAL y recuperación: hacer durable una decisión}

\subsection{El problema que resuelve}

Sin WAL, una alternativa simple sería escribir y sincronizar inmediatamente todas las páginas de tabla e índice modificadas antes de confirmar cada transacción. Eso puede introducir muchas escrituras aleatorias y picos de latencia.

WAL separa dos momentos:

\begin{enumerate}
\item persistir una descripción secuencial suficiente para recuperar la transacción;
\item materializar más tarde las páginas e índices en un orden conveniente.
\end{enumerate}

La garantía de recuperación que usaremos es:

\begin{quote}
Después de recovery deben aparecer todas y solo las transacciones que el sistema confirmó al cliente.
\end{quote}

Esa frase contiene dos obligaciones:

\begin{itemize}
\item \textbf{todas}: ninguna compra confirmada puede desaparecer;
\item \textbf{solo}: una transacción incompleta no puede reaparecer como confirmada.
\end{itemize}

\subsection{Las dos reglas de WAL}

\paragraph{Regla 1: log antes que página.} El registro necesario para REDO debe ser durable antes de que la página modificada pueda llegar al archivo de datos.

Si se invierte el orden, una página podría contener un cambio sin que recovery disponga de evidencia suficiente para interpretarlo o completarlo.

\paragraph{Regla 2: COMMIT antes que ACK.} El registro de COMMIT debe ser durable antes de responder éxito al cliente.

Si se invierte, la aplicación podría mostrar ``compra confirmada'' y, tras un crash, el sistema no podría reconstruir la promesa.

Estas reglas son diferentes. La primera protege la relación entre WAL y páginas. La segunda protege la relación entre el sistema y el cliente.

\subsection{WAL como carril propio}

En un timeline se dibujan carriles para cliente, servicio, memoria del motor, WAL y páginas. WAL no es otro microservicio: es un archivo o conjunto de segmentos administrados por el mismo motor. Se representa por separado para observar tres hechos distintos:

\begin{enumerate}
\item cuándo se agrega un registro al buffer del log;
\item cuándo ese prefijo del log se vuelve durable;
\item cuándo la modificación finalmente aparece en la página de datos.
\end{enumerate}

\subsection{Timeline preciso de CampusShop}

\begin{center}
\begin{tabular}{p{0.08\textwidth}p{0.80\textwidth}}
\toprule
Paso & Evento y estado resultante \\
\midrule
1 & El cliente envía la solicitud 9F2A. Todavía no existe evidencia durable. \\
2 & El motor prepara la fila de \texttt{purchases}, la reducción de \texttt{inventory} y los cambios de índice en memoria. \\
3 & Se agregan registros \texttt{WAL UPDATE}. El append puede seguir en buffers volátiles. \\
4 & Las páginas correspondientes quedan sucias en memoria. \\
5 & El motor agrega un registro \texttt{COMMIT}. Todavía puede no ser durable. \\
6 & El motor hace flush del WAL hasta el LSN del COMMIT. El prefijo relevante queda durable. \\
7 & Recién entonces el servicio envía el ACK o \texttt{200 OK}. \\
8 & Ocurre un crash antes de que todas las páginas hayan sido escritas. \\
9 & Recovery lee el WAL, identifica el COMMIT y ejecuta REDO sobre páginas atrasadas. \\
\bottomrule
\end{tabular}
\end{center}

\begin{uainfo}
\textbf{¿El flush durable comienza en el paso 3?} No. En el paso 3 comienza el append lógico del registro de actualización. El flush ocurre en el paso 6 y vuelve durable un prefijo que llega, como mínimo, hasta el LSN del COMMIT. Como el WAL es secuencial, ese prefijo incluye los registros de actualización anteriores que aún no eran durables.
\end{uainfo}

Formalmente, antes del ACK de una transacción $T$ debe cumplirse:

\[
\texttt{flushLSN} \geq \texttt{commitLSN}(T).
\]

Si parte de los registros anteriores ya había sido persistida por un flush de background o por otro group commit, el dispositivo no necesita reescribir físicamente todo desde el comienzo; basta con avanzar el prefijo durable hasta el COMMIT requerido.

\subsection{Tres puntos de crash y por qué son correctos}

\begin{center}
\begin{tabular}{p{0.28\textwidth}p{0.27\textwidth}p{0.34\textwidth}}
\toprule
Crash & Estado esperado & Razón \\
\midrule
Antes del COMMIT durable & La compra no aparece confirmada & El sistema todavía no emitió un ACK legítimo y recovery no debe presentar trabajo incompleto como confirmado. \\
Después del COMMIT durable, antes de la página & La compra reaparece mediante REDO & La promesa ya fue confirmada y el WAL conserva evidencia suficiente. \\
Después del COMMIT durable y de la página & La compra permanece una sola vez & Recovery debe reconocer que la página ya contiene el cambio o reaplicarlo de forma segura. \\
\bottomrule
\end{tabular}
\end{center}

Los tres resultados obedecen un único criterio: \emph{todas y solo las transacciones confirmadas}.

\subsection{REDO repetible y pageLSN}

Recovery puede sufrir otra caída mientras trabaja. Por eso REDO debe poder repetirse sin duplicar el efecto físico o lógico. Un patrón frecuente consiste en asociar a cada página el LSN del último cambio ya aplicado. Si la página contiene un \texttt{pageLSN} mayor o igual al LSN del registro que se intenta rehacer, el motor sabe que no necesita volver a aplicar ese registro.

Esto no sustituye la idempotencia de negocio. WAL evita perder una historia confirmada del motor; 9F2A evita que dos intentos de la misma compra generen dos compras lógicas.

\subsection{Group commit}

Una sincronización física puede ser costosa. En \textbf{group commit}, varias transacciones agregan sus registros de COMMIT al WAL y comparten un mismo flush durable.

La garantía no se debilita si cada ACK sale después de que el prefijo que contiene su COMMIT se volvió durable. El trade-off es:

\begin{itemize}
\item menor costo promedio por transacción y mayor throughput;
\item una pequeña espera para formar el grupo;
\item posible incremento de latencia de cola bajo carga.
\end{itemize}

\subsection{Checkpoint, RPO y RTO}

Un \textbf{checkpoint} registra progreso suficiente para evitar comenzar recovery desde el origen del log. Simplificadamente, permite distinguir:

\begin{itemize}
\item páginas ya materializadas o cuya escritura fue coordinada;
\item WAL posterior que todavía debe analizarse y reproducirse.
\end{itemize}

Más checkpoints pueden reducir el trabajo de recovery, pero elevan presión de E/S. Checkpoints más espaciados suavizan parte del camino normal, pero aumentan WAL retenido y tiempo de recuperación.

\begin{uawarn}
Un checkpoint no es un backup. El checkpoint acota recuperación local; el backup permite reconstruir datos ante pérdida o corrupción del almacenamiento. Tampoco reemplaza la replicación.
\end{uawarn}

RTO pregunta cuánto puede tardar la recuperación. RPO pregunta cuánto estado confirmado podría perderse bajo el modelo declarado. Con commit sincrónico a WAL local, el RPO frente a crash de proceso puede ser cero, pero el RPO frente a pérdida total del dispositivo requiere otra copia o backup.

\subsection{Caso real: PostgreSQL}

PostgreSQL utiliza WAL para desacoplar la confirmación de una transacción de la escritura inmediata de todas sus páginas. Los registros WAL necesarios se persisten antes de depender de las páginas; al reiniciar, recovery puede hacer roll-forward y rehacer modificaciones confirmadas todavía ausentes de los archivos de datos.

La enseñanza transferible no es un comando de PostgreSQL. Es el orden físico que sostiene la promesa.

\subsection{Cómo se reconstruye la data si las páginas todavía no se materializaron}

El registro \texttt{COMMIT} no contiene por sí solo toda la data de la transacción. Su función es indicar que los registros de modificación anteriores pertenecen a una transacción confirmada. La recuperación combina dos fuentes:

\[
\text{estado recuperado}
=
\text{páginas persistentes anteriores}
+
\operatorname{REDO}(\text{WAL confirmado pendiente}).
\]

Las páginas persistentes aportan un estado base válido. Los registros \texttt{WAL UPDATE} contienen información suficiente para rehacer inserciones, actualizaciones, eliminaciones y cambios de índices que todavía no aparecen en esas páginas. El registro \texttt{COMMIT} aporta la decisión: esos cambios forman parte de la historia confirmada.

Por eso conviene recordar la frase:

\begin{quote}
Los registros de actualización contienen el \emph{cómo}; el registro de COMMIT contiene la \emph{decisión}.
\end{quote}

Si se hubiese perdido todo el dispositivo, el WAL reciente no sería necesariamente una copia completa de la base. Ese escenario requiere un backup base más una secuencia continua de WAL o una réplica adecuada. WAL de crash recovery y backup resuelven fallas diferentes.

\section{MVCC: acceso concurrente a una historia de datos}

\subsection{Problema antes del mecanismo}

Un reporte del comedor comienza a las 12:00 y tarda 40 segundos. Durante ese intervalo entran 900 compras. Hay tres posibilidades ingenuas:

\begin{enumerate}
\item bloquear todas las escrituras hasta que termine el reporte;
\item permitir que el reporte mezcle filas de momentos diferentes;
\item conservar versiones y aplicar una regla de visibilidad coherente.
\end{enumerate}

MVCC implementa la tercera alternativa.

\subsection{Definición y límite de la metáfora}

\textbf{Control de concurrencia multiversión} significa que una escritura crea una versión nueva y el motor conserva temporalmente versiones anteriores. Cada transacción lectora usa un snapshot para elegir qué versión pertenece a la historia que puede observar.

La metáfora de una fotografía es útil, pero tiene límites:

\begin{itemize}
\item no se copia físicamente toda la base al comenzar la transacción;
\item el snapshot es una regla de visibilidad;
\item una fotografía coherente no implica por sí sola serializabilidad;
\item conservar versiones consume espacio y exige limpieza.
\end{itemize}

\subsection{Timeline de visibilidad}

Suponga que $v_1=120$ está confirmada.

\begin{enumerate}
\item El reporte $T_R$ comienza y fija el snapshot $S_R$.
\item $T_R$ lee $v_1$.
\item Una transacción escritora $T_W$ crea $v_2=135$ sin borrar inmediatamente $v_1$.
\item $T_W$ ejecuta COMMIT.
\item Una segunda lectura dentro de $T_R$ sigue observando $v_1$ si el nivel de aislamiento conserva el snapshot.
\item Una transacción nueva puede observar $v_2$.
\item Cuando termina $T_R$ y ningún snapshot necesita $v_1$, la versión anterior puede ser reclamable.
\end{enumerate}

Una formulación conceptual de visibilidad es:

\begin{quote}
Una versión es visible si la transacción que la creó pertenece a la historia confirmada que el snapshot puede observar y su reemplazo o eliminación todavía no era visible para ese snapshot.
\end{quote}

Los motores concretos implementan esta regla con identificadores de transacción, timestamps, listas de transacciones activas u otros metadatos.

\subsection{Historia retenida, VACUUM y bloat}

Una versión anterior no puede eliminarse mientras un snapshot activo pueda necesitarla. Una transacción muy larga puede impedir la reclamación de muchas versiones.

En PostgreSQL, \texttt{VACUUM} recupera espacio reutilizable cuando esas versiones dejan de ser visibles. Si la limpieza se retrasa aparecen:

\begin{itemize}
\item más páginas que leer;
\item mayor presión de caché;
\item índices con referencias a versiones obsoletas;
\item mayor E/S;
\item crecimiento de almacenamiento o \emph{bloat}.
\end{itemize}

MVCC no elimina el costo de concurrencia: cambia bloqueo por administración de historia.

\subsection{MVCC, aislamiento y serializabilidad}

MVCC es un mecanismo, no un nivel de aislamiento único. Un motor multiversión puede ofrecer Read Committed, Repeatable Read, Snapshot Isolation o Serializable según sus reglas de snapshot y detección de conflictos.

Considere el invariante: ``al menos una de las dos guardias debe permanecer activa''.

\begin{itemize}
\item $T_1$ lee $A=1,B=1$ y escribe $A=0$.
\item $T_2$ lee $A=1,B=1$ y escribe $B=0$.
\item Como escriben filas distintas, ambas pueden confirmar bajo Snapshot Isolation.
\item El resultado $A=0,B=0$ viola el invariante.
\end{itemize}

Este fenómeno es \textbf{write skew}. MVCC mantuvo snapshots coherentes, pero no garantizó serializabilidad. La prevención puede requerir serializable isolation, locks explícitos, materializar el conflicto en una misma fila o rediseñar el invariante.

\section{Índices y EXPLAIN: medir antes de optimizar}

\subsection{Transición desde MVCC: visibilidad no es localización}

MVCC responde una pregunta lógica: \emph{entre todas las versiones posibles, cuál puede observar esta transacción según su snapshot}. Sin embargo, antes de aplicar esa regla el motor debe localizar filas y versiones candidatas.

Una estructura de acceso responde una pregunta física distinta: \emph{cómo llegar a un subconjunto pequeño sin recorrer todas las páginas}. Por eso el flujo conceptual es:

\[
\text{consulta SQL}
\rightarrow
\text{planificador}
\rightarrow
\text{ruta física}
\rightarrow
\text{filas/versiones candidatas}
\rightarrow
\text{reglas MVCC}.
\]

El índice no decide qué versión es correcta. Reduce el espacio de búsqueda; MVCC conserva la semántica de la lectura.

\subsection{Árbol B como estructura adicional}

Un \textbf{árbol B} o \textbf{B-tree} es un árbol ordenado, balanceado y de múltiples hijos. Sus nodos suelen corresponder a páginas de almacenamiento. Los nodos internos conservan claves separadoras y referencias a otras páginas; las hojas conservan claves ordenadas y referencias a filas o páginas de datos.

El \textbf{fan-out} es la cantidad de claves y referencias que caben en una página del índice. Como suele ser alto, un árbol con millones de entradas puede mantener poca altura y requerir pocas lecturas para llegar desde la raíz hasta una hoja.

Esta estructura favorece:

\begin{itemize}
\item búsquedas por igualdad;
\item rangos y recorridos ordenados;
\item combinación de predicados mediante índices compuestos;
\item detención temprana cuando una consulta usa \texttt{LIMIT}.
\end{itemize}

El árbol no reemplaza la tabla ni almacena necesariamente la fila completa. Es una estructura redundante que ocupa espacio y debe mantenerse cuando cambian los datos.

\subsection{Consulta de CampusShop}

\begin{verbatim}
SELECT id, total, confirmed_at
FROM purchases
WHERE student_id = 8421
  AND status = 'CONFIRMED'
ORDER BY confirmed_at DESC
LIMIT 20;
\end{verbatim}

La pregunta contiene tres propiedades importantes:

\begin{enumerate}
\item igualdad por \texttt{student\_id};
\item igualdad por \texttt{status};
\item orden descendente por \texttt{confirmed\_at} con corte temprano.
\end{enumerate}

Una hipótesis de índice es:

\begin{verbatim}
CREATE INDEX idx_purchase_student_status_date
ON purchases(student_id, status, confirmed_at DESC);
\end{verbatim}

El orden de columnas refleja el patrón de acceso: primero restringe por estudiante, después por estado y finalmente recorre las fechas ya ordenadas de manera descendente. Si las primeras veinte entradas cumplen la visibilidad MVCC, el motor puede detener el recorrido temprano.

Una ruta conceptual es:

\[
(8421,\texttt{CONFIRMED})
\rightarrow
\text{raíz}
\rightarrow
\text{página hoja}
\rightarrow
\text{fechas descendentes}
\rightarrow
20\text{ resultados}.
\]

\texttt{CREATE INDEX} es una sentencia DDL porque modifica la estructura del esquema.

\subsection{EXPLAIN frente a EXPLAIN ANALYZE}

\texttt{EXPLAIN} muestra el plan elegido y estimaciones de costo y cardinalidad. \texttt{EXPLAIN ANALYZE} ejecuta la consulta y agrega mediciones reales. Con \texttt{BUFFERS} puede mostrar bloques leídos o encontrados en caché.

Una regla práctica es leer el plan de abajo hacia arriba:

\[
\text{acceso físico} \rightarrow \text{filtro} \rightarrow \text{orden} \rightarrow \text{límite}.
\]

En el ejemplo didáctico sin índice aparecen:

\begin{itemize}
\item \texttt{Seq Scan};
\item miles de filas descartadas;
\item muchos buffers tocados;
\item ordenamiento separado;
\item veinte filas finales.
\end{itemize}

Con el índice compuesto pueden aparecer:

\begin{itemize}
\item \texttt{Index Scan};
\item pocas filas visitadas;
\item menos buffers;
\item orden ya compatible con la consulta;
\item detención temprana por \texttt{LIMIT}.
\end{itemize}

Las cifras de las diapositivas son ilustrativas. El resultado real depende de tamaño, distribución, estadísticas, caché, hardware y versión del motor.

\subsection{Por qué un índice puede no utilizarse}

El planificador puede preferir un scan secuencial cuando:

\begin{itemize}
\item la consulta necesita una fracción grande de la tabla;
\item la tabla es pequeña;
\item la selectividad es baja;
\item las estadísticas están desactualizadas;
\item el acceso aleatorio al índice y a la tabla sería más caro;
\item otro índice o plan ofrece menor costo estimado.
\end{itemize}

\subsection{Costo de sobreindexar}

Cada índice puede agregar:

\begin{itemize}
\item escrituras adicionales;
\item más WAL;
\item más páginas y espacio;
\item presión de caché;
\item mantenimiento durante VACUUM o compactación;
\item splits o reorganización;
\item mayor duración de cargas y restauraciones.
\end{itemize}

La decisión correcta no es ``crear índices''. Es formular una hipótesis de acceso, observar el plan y medir el workload completo.

\section{LSM-tree: escrituras rápidas y organización posterior}

\subsection{Motivación}

En una carga de telemetría del campus, sensores de aulas, energía y clima pueden producir decenas de miles de eventos por segundo. Actualizar muchas páginas e índices de manera aleatoria puede ser costoso. Una LSM-tree desplaza gran parte de ese trabajo:

\begin{enumerate}
\item registra la escritura en WAL;
\item la inserta en una memtable ordenada en memoria;
\item congela la memtable cuando alcanza un umbral;
\item la vuelca como SSTable inmutable en nivel L0;
\item fusiona SSTables posteriormente mediante compactación.
\end{enumerate}

No elimina el costo de escribir: transforma el lugar y el momento en que se paga.

\subsection{Vocabulario LSM}

\begin{description}
\item[Memtable.] Tabla ordenada en memoria que recibe escrituras recientes.
\item[Memtable inmutable.] Memtable cerrada a nuevas escrituras y preparada para volcarse.
\item[SSTable --- Sorted String Table.] Archivo ordenado e inmutable de pares clave--valor.
\item[Run.] Secuencia o archivo ya ordenado que puede participar en una fusión.
\item[Nivel L0, L1, L2.] Conjuntos de SSTables organizados por políticas de tamaño y solapamiento.
\item[Tombstone.] Marca de borrado que representa una eliminación y debe conservarse hasta que sea seguro olvidar versiones anteriores.
\item[Merge ordenado.] Fusión que recorre runs ya ordenados, selecciona la menor clave disponible y produce otra secuencia ordenada sin ordenar todos los datos desde cero.
\item[Compactación.] Merge ordenado que además reduce solapamiento y elimina versiones obsoletas cuando es seguro.
\item[Backlog.] Cantidad de bytes o archivos todavía pendientes de compactación.
\item[Write stall.] Pausa o desaceleración temporal de escrituras para evitar que el backlog vuelva inestable al motor.
\end{description}

\subsection{Ruta de escritura}

El camino crítico típico es:

\[
(k,v) \rightarrow WAL \rightarrow memtable \rightarrow memtable\ inmutable
\rightarrow flush \rightarrow SSTable\ L0.
\]

L0 puede contener varios archivos con rangos solapados. La compactación reorganiza esos archivos hacia niveles más estables.

\subsection{Merge ordenado y tombstones}

Suponga:

\begin{itemize}
\item SSTable A: $a,c_1,e,g$;
\item SSTable B: $b,c_2,f,h$;
\item tombstone: eliminar $c$.
\end{itemize}

El merge ordenado compara las claves de entrada, conserva la versión vigente y produce una nueva SSTable. La marca de borrado no puede eliminarse inmediatamente si existe el riesgo de que una versión antigua de $c$ sobreviva en otro nivel; de lo contrario, esa versión podría ``resucitar'' después de una compactación parcial.

\subsection{Filtro de Bloom}

Un \textbf{filtro de Bloom} es un resumen probabilístico y compacto de las claves de una SSTable. Para una búsqueda exacta no devuelve el valor. Responde:

\begin{itemize}
\item \textbf{definitivamente ausente}: al menos uno de los bits consultados es cero; esa SSTable puede omitirse para esa clave;
\item \textbf{posiblemente presente}: todos los bits consultados son uno; el motor debe consultar el índice y quizá el bloque de datos.
\end{itemize}

Puede producir falsos positivos: decir ``posiblemente presente'' cuando la clave no está. No debe producir falsos negativos: nunca puede descartar una SSTable que realmente contiene la clave.

La expresión ``archivo imposible'' es ambigua y debe evitarse. La formulación correcta es:

\begin{quote}
SSTable que puede omitirse para esta búsqueda porque el filtro demuestra que la clave no está allí.
\end{quote}

El filtro no reemplaza el índice del archivo, la caché de bloques, la elección de la versión visible ni las reglas sobre tombstones.

\subsection{Ruta de lectura}

Una búsqueda puntual puede seguir este orden:

\begin{enumerate}
\item revisar memtables, que contienen lo más reciente;
\item seleccionar SSTables candidatas por nivel y rango;
\item consultar el filtro de Bloom;
\item omitir archivos con ausencia definitiva;
\item consultar el índice de los archivos con presencia posible;
\item usar la caché o leer el bloque del dispositivo;
\item reconciliar versiones y respetar tombstones.
\end{enumerate}

\section{Compactación, amplificaciones y pausas de escritura}

\subsection{Tres amplificaciones}

\paragraph{Write amplification.} Relación entre bytes físicos escritos y bytes lógicos ingresados:

\[
WA = \frac{\text{bytes físicos escritos}}{\text{bytes lógicos ingresados}}.
\]

WAL, índices, replicación y compactación pueden aumentar el numerador.

\paragraph{Read amplification.} Cantidad de estructuras, archivos o bloques que deben examinarse para responder una lectura.

\paragraph{Space amplification.} Relación entre espacio físico ocupado y datos lógicos vivos. Versiones, tombstones, archivos solapados y espacio temporal contribuyen.

Estas dimensiones no se minimizan simultáneamente. Cada política redistribuye costo.

\subsection{Backlog y capacidad}

Si el sistema recibe 90 MB/s y la compactación procesa 60 MB/s, el déficit sostenido es:

\[
90 - 60 = 30\ \text{MB/s}.
\]

Mientras esa desigualdad continúe, el backlog crece. La cadena causal es:

\[
\begin{aligned}
\text{ingreso} > \text{capacidad}
&\Rightarrow \text{backlog creciente} \\
&\Rightarrow \text{más archivos L0} \\
&\Rightarrow \text{más trabajo de lectura y espacio} \\
&\Rightarrow \text{latencia alta} \\
&\Rightarrow \text{slowdown o stall}.
\end{aligned}
\]

\subsection{Pausa de escritura como protección}

Una pausa de escritura es backpressure interno. Reduce o detiene temporalmente nuevas escrituras para impedir que una cola interna crezca sin límite.

No crea capacidad. Si ocurre frecuentemente o dura demasiado, puede indicar:

\begin{itemize}
\item carga sostenida mayor que la capacidad;
\item política de compactación inadecuada;
\item disco insuficiente o saturado;
\item demasiados archivos L0;
\item claves o particiones calientes;
\item competencia con lecturas, snapshots o backups.
\end{itemize}

\subsection{Métricas que deben leerse juntas}

\begin{itemize}
\item bytes lógicos ingresados por segundo;
\item bytes leídos y escritos por compactación;
\item write amplification;
\item backlog en bytes;
\item cantidad de archivos L0;
\item tiempo en slowdown y stall;
\item p95 y p99 de escritura;
\item uso de disco, CPU y caché;
\item throughput efectivo admitido.
\end{itemize}

Ninguna cifra aislada basta. Backlog creciente junto con p99 creciente y tiempo de stall revela una cadena causal más fuerte que cualquiera de esas señales por separado.

\subsection{Leveled frente a tiered}

\begin{center}
\begin{tabular}{p{0.22\textwidth}p{0.33\textwidth}p{0.33\textwidth}}
\toprule
Dimensión & Leveled & Tiered / universal \\
\midrule
Lecturas & Menos archivos solapados & Más runs potenciales \\
Escrituras & Más reescritura en merges & Menos reescritura temprana \\
Espacio & Mejor control & Más espacio temporal \\
Uso típico & Búsquedas puntuales sensibles & Ingestión intensa \\
\bottomrule
\end{tabular}
\end{center}

La elección depende del patrón de acceso, hardware, retención, amplificaciones y SLO.

\section{B-tree y LSM: dos hipótesis, no una ley}

\begin{center}
\begin{tabular}{p{0.20\textwidth}p{0.34\textwidth}p{0.34\textwidth}}
\toprule
Criterio & Árbol B y páginas & LSM y archivos ordenados \\
\midrule
Camino de escritura & Actualiza páginas e índices & WAL, memtable, flush y compactación \\
Búsqueda puntual & Pocas páginas si el índice ayuda & Puede revisar varios archivos; Bloom y caché reducen trabajo \\
Rangos & Natural por orden del árbol & Bueno si la compactación limita runs \\
Costo diferido & Splits, VACUUM, bloat y reconstrucción & Tombstones, compactación y stalls \\
Carga frecuente & OLTP equilibrado o read-heavy & Ingestión sostenida o write-heavy \\
Pregunta clave & ¿Qué consulta justifica el índice? & ¿Puede la compactación sostener la tasa? \\
\bottomrule
\end{tabular}
\end{center}

\paragraph{CampusShop.} Transacciones cortas, inventario, idempotencia y consultas por estudiante y fecha. Una hipótesis razonable es WAL + MVCC + árbol B.

\paragraph{Telemetría del campus.} Ingestión masiva, claves por dispositivo y tiempo, retención temporal y vistas eventualmente frescas. Una hipótesis razonable es WAL + LSM + compactación orientada a tiempo.

No basta con elegir. Hay que medir percentiles, throughput sostenible, RPO, RTO, amplificaciones, backlog, hit rate y espacio.

\section{Método reusable de diseño}

Para cualquier subsistema de almacenamiento:

\begin{enumerate}
\item \textbf{Carga}: proporción de lecturas/escrituras, tamaño, sesgo, crecimiento, rangos y retención.
\item \textbf{Garantía}: qué significa durable, visible y correcto bajo una falla concreta.
\item \textbf{Mecanismo}: WAL, MVCC, índice, árbol B, LSM, replicación u otra combinación.
\item \textbf{Costo}: latencia, E/S, espacio, complejidad, mantenimiento y recuperación.
\item \textbf{Evidencia}: plan, métrica y experimento capaz de refutar la elección.
\end{enumerate}

Aplicado al marco del curso:

\[
D=(P,F,G,S,T).
\]

Para telemetría:

\begin{itemize}
\item $P$: ingerir y consultar eventos a alta tasa;
\item $F$: crash, pérdida de nodo, backlog, picos y disco lento;
\item $G$: RPO acotado, RTO definido y rangos consultables;
\item $S$: WAL, LSM, replicación, compactación y retención;
\item $T$: write amplification, frescura eventual, espacio temporal y stalls.
\end{itemize}

\section{Conexión con el Trabajo Final Integrador}

La estrategia de datos del Trabajo Final no debe limitarse a nombrar un producto. Cada grupo debería documentar:

\begin{itemize}
\item operación crítica e identidad lógica;
\item patrón de lectura, escritura, crecimiento y sesgo;
\item definición concreta de durabilidad y modelo de falla;
\item frontera de confirmación y estrategia de recuperación;
\item nivel de aislamiento y anomalías que no se toleran;
\item ruta física de las consultas críticas e índices propuestos;
\item estructura de almacenamiento y política de mantenimiento;
\item SLO, RPO y RTO;
\item métricas y experimento capaz de refutar la elección.
\end{itemize}

La defensa debe poder responder: \emph{qué problema resuelve la elección, qué garantía sostiene, qué costo desplaza y bajo qué evidencia se revisaría}.

\section{Errores conceptuales frecuentes}

\begin{itemize}
\item confundir \texttt{write()} exitoso con dato durable;
\item asumir que replicar en RAM sustituye persistencia;
\item decir que el flush durable ``empieza'' cuando se agrega el primer registro al WAL;
\item creer que COMMIT exige haber escrito todas las páginas;
\item atribuir exactamente-una-vez a WAL sin idempotencia de negocio;
\item usar snapshot como sinónimo de copia física completa;
\item concluir que MVCC implica serializabilidad;
\item crear índices sin medir su costo de escritura;
\item leer \texttt{EXPLAIN} como una promesa universal de tiempo;
\item usar ``archivo imposible'' en vez de ``SSTable que puede omitirse para esta búsqueda'';
\item creer que un filtro de Bloom devuelve valores o elimina falsos positivos;
\item tratar compactación como mantenimiento gratuito;
\item interpretar un stall como capacidad adicional;
\item elegir un motor por reputación y no por carga y evidencia.
\end{itemize}

\section{Glosario operativo compacto}

\begin{description}
\item[Durable.] Reconstruible después de la falla declarada.
\item[COMMIT durable.] Registro de confirmación incluido en el prefijo durable del WAL.
\item[Flush durable.] Volcado que hizo persistentes los bytes requeridos.
\item[WAL.] Registro anticipado de escritura para recuperación.
\item[LSN.] Posición ordenada dentro del WAL.
\item[Snapshot.] Regla de visibilidad de una lectura.
\item[Serializabilidad.] Equivalencia con alguna ejecución secuencial.
\item[DDL.] Sentencias que definen estructuras del esquema.
\item[SSTable.] Archivo ordenado e inmutable.
\item[Merge ordenado.] Fusión de secuencias previamente ordenadas.
\item[Filtro de Bloom.] Resumen probabilístico que permite descartar ausencias definitivas.
\item[Backlog.] Trabajo físico pendiente de mantenimiento.
\item[Write stall.] Backpressure aplicado por el motor a nuevas escrituras.
\end{description}

\section{Autoevaluación}

\begin{enumerate}
\item ¿Qué diferencia existe entre una operación aceptada, escrita y durable?
\item ¿Por qué WAL aparece como un carril separado aunque sea parte del mismo motor?
\item ¿Cuáles son las dos reglas de WAL y qué falla evita cada una?
\item ¿Qué prefijo debe volver durable el flush antes del ACK?
\item ¿Por qué el flush durable no ``comienza'' con el append del \texttt{WAL UPDATE}?
\item ¿Por qué una página puede seguir en memoria después de un COMMIT correcto?
\item ¿Cómo hace recovery para distinguir una transacción confirmada de una incompleta?
\item ¿Por qué REDO debe ser repetible?
\item ¿Qué diferencia hay entre checkpoint, backup y réplica?
\item ¿Qué es un snapshot sin usar la idea de copiar toda la base?
\item ¿Cuándo puede eliminarse una versión antigua?
\item ¿Por qué una transacción larga puede aumentar bloat?
\item ¿Por qué MVCC no implica serializabilidad?
\item ¿Qué muestra \texttt{EXPLAIN} y qué agrega \texttt{EXPLAIN ANALYZE}?
\item ¿Por qué un scan secuencial puede ser correcto aunque exista un índice?
\item ¿Por qué importa el orden de columnas en un índice compuesto?
\item ¿Qué pasos forman el camino de escritura de una LSM-tree?
\item ¿Qué es un merge ordenado?
\item ¿Qué garantiza y qué no un filtro de Bloom?
\item ¿Qué señales muestran que la compactación no sostiene la tasa de ingreso?
\end{enumerate}

\section{Bibliografía guiada}

\begin{itemize}
\item Martin Kleppmann, \emph{Designing Data-Intensive Applications}, capítulos 3, 7 y 9: estructuras de almacenamiento, transacciones, aislamiento y consistencia.
\item PostgreSQL Documentation, \emph{Write-Ahead Logging}: relación entre WAL, páginas y crash recovery.
\item PostgreSQL Documentation, \emph{Routine Vacuuming} y \emph{MVCC}: versiones, visibilidad y limpieza.
\item O'Neil et al., \emph{The Log-Structured Merge-Tree}: fundamento académico de la estructura LSM.
\item RocksDB Wiki, \emph{Compaction}, \emph{Compaction Stats} y \emph{Write Stalls}: métricas y protección de capacidad.
\item Jim Gray y Andreas Reuter, \emph{Transaction Processing}: fundamentos de logging y recuperación.
\end{itemize}

\begin{uakeybox}
\textbf{Síntesis final.} WAL hace recuperable una decisión confirmada; MVCC organiza una historia accedida concurrentemente; índices y \texttt{EXPLAIN} permiten justificar rutas físicas; una LSM-tree desplaza organización al mantenimiento; compactación y stalls controlan ese costo diferido. Ningún mecanismo elimina el costo: decide cuándo, dónde y cómo se paga.
\end{uakeybox}

\end{document}
